\documentclass[10pt,a4paper]{amsart}
\usepackage[margin=19mm]{geometry}
\usepackage{amsmath,amssymb,mathtools}
\usepackage[scr=rsfs,cal=euler]{mathalfa}
\usepackage{xfrac}
\usepackage[unicode,hidelinks,bookmarks=false]{hyperref}
\newcommand{\ie}{i.e.\ }
\newcommand{\cf}{cf.\ }
\newcommand{\resp}{respectively\ }
\newcommand{\Subsection}{\subsection}
\providecommand{\nobreakdash}{\nobreak\hbox{-}}
\setlength{\emergencystretch}{2em}
\multlinegap0pt
\usepackage{amssymb}
\usepackage{mathtools}
\usepackage{enumerate}
\usepackage{xcolor}
\newtheorem{lemma}{Lemma}
 \newcommand{\M}{\mathfrak{M}}
 \newcommand{\m}{\mathfrak{m} }
\title{Bounds on the second Hilbert coefficient and the depth of the associated graded ring}
\author{Clare D'Cruz, Mousumi Mandal, Shruti Priya}
\date{Source: arXiv:2607.29311v1 (CC BY 4.0)}
\begin{document}
\maketitle

\noindent\emph{Source and licence.} Bounds on the second Hilbert coefficient and the depth of the associated graded ring. Version arXiv:2607.29311v1 (CC BY 4.0), distributed by arXiv under the Creative Commons Attribution 4.0 International licence, http://creativecommons.org/licenses/by/4.0/, as recorded in the arXiv licence metadata for this exact version and shown on the abstract page https://arxiv.org/abs/2607.29311v1. Full licence notice: \texttt{licenses/arxiv-2607.29311-CC-BY-4.0.txt}.\par\smallskip

\noindent\textit{Editorial scope.} Counted unit: the article's lemma comparing (source0.tex lines 1730--1736). The article's complete original proof of it is delivered with it (source0.tex lines 1738--1758, one \texttt{proof} environment of 21 lines). No mathematics is written, repaired or re-derived: the statement and the proof are the article's own lines, in the article's own notation, and no retained line is substituted or re-flowed.  Retained uncounted as the source-local context the proof needs: lines 1718--1723.  Nothing was omitted from any retained span, so the package carries no omission record.  No retained line contains a citation, so the wrapper carries no bibliography and no bibliography entry.  No retained line contains a figure environment, an image-inclusion command or drawing code, so no image is delivered and the package declares no asset dependency.  Editorial formatting only, in addition: an AMS \texttt{amsart} preamble that restates the article's own declarations and macros used by the retained lines, namely the environments \texttt{lemma} on separate counters, together with the standard packages those lines need.  The retained environments print in this document as Lemma 1 in order of appearance, whereas the article prints the same items with the numbering of its own full document.  The complete native source of the article as distributed in the arXiv e-print for this exact version is bundled unchanged as \texttt{sources/206471/source0.tex}.
\begin{equation} \label{4}
    \begin{aligned}
        0 \longrightarrow \mathcal{B}(x,R)  &  \longrightarrow H^{0}_{\M}\left(L^{I}(R)\right)(-1) \longrightarrow H^{0}_{\M}\left(L^{I}(R)\right) \longrightarrow H^{0}_{\M}\left(L^{I'}(R')\right) \\
     & \xlongrightarrow{\delta} H^{1}_{\M}\left(L^{I}(R)\right)(-1) \longrightarrow H^{1}_{\M}\left(L^{I}(R)\right) \longrightarrow H^{1}_{\M}\left(L^{I'}(R')\right) \cdots.
    \end{aligned}
\end{equation}

\begin{lemma} \label{comparing}
Let $(R, \m)$ be a Cohen-Macaulay local ring and $I$ be an $\m$-primary ideal. 
 Then 
 \begin{equation*}
     \lambda(\mathcal{B}(x,R))\leq \lambda\left(H^{0}_{\M}\left(L^{I'}(R')\right)\right).
 \end{equation*}
\end{lemma}

\begin{proof}
    From the long exact sequence of local cohomology modules in (\ref{4}), we have the following exact sequence:
    \begin{equation} \label{22}
       0 \longrightarrow \mathcal{B}(x,R)   \longrightarrow H^{0}_{\M}\left(L^{I}(R)\right)(-1) 
        \longrightarrow H^{0}_{\M}\left(L^{I}(R)\right) 
        \longrightarrow H^{0}_{\M}\left(L^{I'}(R')\right) 
        \longrightarrow C 
        \longrightarrow 0,
    \end{equation}

    where $C:=\textrm{image} (\delta)$. 
Taking  lengths of modules in (\ref{22}), we get
\begin{align*}
       \lambda\left(\mathcal{B}(x,R)\right)
 & = \lambda\left(H^{0}_{\M}\left(L^{I}(R)\right)(-1)\right)
  -\lambda\left(H^{0}_{\M}\left(L^{I}(R)\right)\right)
  +\lambda\left(H^{0}_{\M}\left(L^{I'}(R')\right)\right)
  - \lambda(C)\\
 & \leq \lambda\left(H^{0}_{\M}\left(L^{I'}(R')\right)\right)  \text{ (as  $\lambda\left(H^{0}_{\M}\left(L^{I}(R)\right)\right) < \infty)$}. \quad \quad\qedhere
\end{align*} 
\end{proof}

\end{document}
